\documentclass[11pt,a4paper]{article}

\usepackage{fontspec}
\setmainfont{DejaVu Serif}[
  BoldFont={DejaVu Serif Bold},
  ItalicFont={DejaVu Serif},
  ItalicFeatures={FakeSlant=0.18},
  BoldItalicFont={DejaVu Serif Bold},
  BoldItalicFeatures={FakeSlant=0.18}
]
\setsansfont{DejaVu Sans}
\usepackage{polyglossia}
\setmainlanguage{english}
\usepackage[a4paper,margin=20.5mm]{geometry}
\usepackage{microtype}
\usepackage{setspace}
\setstretch{1.012}
\usepackage{amsmath,amssymb,amsthm,mathtools,bm}
\usepackage{booktabs,array,tabularx}
\usepackage{enumitem}
\usepackage{xcolor}
\usepackage{tikz}
\usetikzlibrary{arrows.meta,positioning,calc}
\usepackage[unicode,hidelinks]{hyperref}
\usepackage[nameinlink,noabbrev]{cleveref}
\emergencystretch=3em
\allowdisplaybreaks

\hypersetup{
  pdftitle={Regular Polyhedra under Typed Observation Filters},
  pdfauthor={Stas, Independent Research Program},
  pdfsubject={Abstract polyhedra, geometric realizations, flag transitivity, observation-filter preorders, Petrials, and finite-resolution identification},
  pdfkeywords={regular polyhedra, abstract polytopes, flag transitivity, string C-groups, Petrie duality, observation geometry, identifiability}
}

\definecolor{obsblue}{RGB}{37,83,138}
\definecolor{obsred}{RGB}{150,45,45}
\definecolor{obsgreen}{RGB}{30,115,78}
\definecolor{obsviolet}{RGB}{94,64,145}
\definecolor{lightblue}{RGB}{239,246,253}
\definecolor{lightred}{RGB}{253,242,242}
\definecolor{lightgreen}{RGB}{239,249,244}
\definecolor{lightviolet}{RGB}{246,242,252}

\theoremstyle{definition}
\newtheorem{definition}{Definition}[section]
\newtheorem{model}[definition]{Model}
\theoremstyle{plain}
\newtheorem{theorem}[definition]{Theorem}
\newtheorem{proposition}[definition]{Proposition}
\newtheorem{lemma}[definition]{Lemma}
\newtheorem{corollary}[definition]{Corollary}
\theoremstyle{remark}
\newtheorem{remark}[definition]{Remark}
\newtheorem{warning}[definition]{Boundary}

\newcommand{\R}{\mathbb{R}}
\newcommand{\Aut}{\operatorname{Aut}}
\newcommand{\Isom}{\operatorname{Isom}}
\newcommand{\Flag}{\operatorname{Flag}}
\newcommand{\Ob}{\operatorname{Ob}}
\newcommand{\Ker}{\operatorname{Ker}}
\newcommand{\ord}{\operatorname{ord}}
\newcommand{\diam}{\operatorname{diam}}
\newcommand{\norm}[1]{\left\lVert #1\right\rVert}
\newcommand{\abs}[1]{\left\lvert #1\right\rvert}
\newcommand{\Ccal}{\mathcal C}
\newcommand{\Pcal}{\mathcal P}
\newcommand{\Ocal}{\mathcal O}
\newcommand{\Qcal}{\mathcal Q}
\newcommand{\Fcal}{\mathcal F}
\newcommand{\Gcal}{\mathcal G}
\newcolumntype{L}[1]{>{\raggedright\arraybackslash}p{#1}}
\newcolumntype{Y}{>{\raggedright\arraybackslash}X}

\title{\bfseries Regular Polyhedra under Typed Observation Filters\\
\large Abstract incidence, realization symmetry, Petrials, and identifiability}
\author{Stas\\
\small Independent Research Program; ORCID: 0009-0000-2294-705X}
\date{Strict GO Core reference revision v1.1 --- 28 July 2026}

\begin{document}
\maketitle

\begin{center}
\fcolorbox{obsblue}{lightblue}{%
\parbox{0.92\linewidth}{\small
\textbf{Status.} This document supersedes the bilingual working note
v1.0 and is an application-layer reference migrated to GO Core v0.2.
It does not introduce a new polyhedron, classification, realization
theorem, or symmetry algorithm. Its contribution is a typed separation
of abstract incidence, geometric realization, admissibility,
information loss, exact symmetry, and finite-resolution inference.}}
\end{center}

\begin{abstract}
The phrase ``regular polyhedron'' changes meaning when convex
polyhedra, star polyhedra, abstract maps, or discrete skeletal
realizations are placed in the same discussion.  This reference fixes
the object class before counting.  An abstract polyhedron is a ranked
poset satisfying the diamond condition and strong flag connectivity;
abstract regularity means transitivity of \(\Aut(\Pcal)\) on flags.
A geometric realization \(\beta\) has a generally smaller Euclidean
symmetry group \(G_\beta\), so an abstractly regular object need not be
geometrically regular in a chosen realization.  Observation
specifications are defined on full subgroupoids of realized objects.
Admissibility enlargement and informational coarsening are two
different preorders: \(\Ccal_1\subseteq\Ccal_2\) enlarges the admitted
domain, whereas \(Q_{\rm coarse}=\kappa\circ Q_{\rm fine}\) loses
distinctions.  They cannot be compressed into an undefined inclusion
of tuples.  A cube and its Petrial give an exact non-injectivity
control: the vertex-edge skeleton is unchanged, while six square faces
are replaced by four Petrie hexagons and the Euler characteristic
changes from \(2\) to \(0\).  Under bounded noise, approximate
symmetry residuals and a half-separation criterion give conditional
identifiability, not proof of exact symmetry.  Counts such as five
Platonic solids or \(48=18+30\) discrete regular skeletal polyhedra in
\(\R^3\) are therefore convention-scoped results, not values of one
universal observation-independent enumeration.
\end{abstract}

\tableofcontents

\section{Scope and the category of objects}

\begin{definition}[Abstract polyhedron]
A rank-\(3\) abstract polyhedron \(\Pcal\) is a partially ordered set
with a rank map to \(\{-1,0,1,2,3\}\), unique least and greatest
faces, and the following properties.
\begin{enumerate}[label=(\roman*),leftmargin=1.6em]
  \item Every maximal chain, or \emph{flag}, contains exactly one face
  of each rank.
  \item If \(F<G\) have ranks \(i-1\) and \(i+1\), exactly two
  rank-\(i\) faces lie strictly between them (diamond condition).
  \item Every section is flag-connected, equivalently the usual
  strong flag-connectivity condition holds.
\end{enumerate}
The proper faces of ranks \(0,1,2\) are vertices, edges, and faces.
A reduced flag is the incident triple \((v,e,f)\).
\end{definition}

The tuple \((V,E,F,I)\) from the legacy note is only bookkeeping; by
itself it does not impose thinness, connectedness, rank, or section
axioms.  Those axioms are required before standard results about
polyhedra and flags are used.

\begin{definition}[Realized-polyhedron groupoid]
Let \(\mathbf{RPol}_3\) be the groupoid whose objects are
\[
X=(\Pcal,\beta,\mathsf T),
\]
where \(\Pcal\) is an abstract polyhedron, \(\beta\) is a declared
realization in \(\R^3\), and \(\mathsf T\) records the realization
convention: straight or other edges, planar, skew, star, zig-zag, or
helical faces, discreteness, local finiteness, and faithfulness.
An arrow \(X\to X'\) is a rank-preserving incidence isomorphism
\(\varphi:\Pcal\to\Pcal'\) together with an ambient isometry \(a\)
intertwining the realizations and all declared edge and face data.
\end{definition}

The convention \(\mathsf T\) is part of the type.  A convex
face-lattice realization, a regular map, and a skeletal apeirohedron
are not silently treated as objects of one unqualified category.

\section{Flags, abstract regularity, and string C-groups}

By the diamond condition, every flag \(\Phi\) has a unique
\(i\)-adjacent flag \(\Phi^i\), differing only in its rank-\(i\)
face, for \(i=0,1,2\).  The resulting edge-coloured flag graph is
connected.

\begin{definition}[Abstract regularity]
The automorphism group \(\Gamma(\Pcal)=\Aut(\Pcal)\) consists of all
rank-preserving order automorphisms.  The abstract polyhedron is
regular when \(\Gamma(\Pcal)\) acts transitively on
\(\Flag(\Pcal)\).
\end{definition}

\begin{proposition}[Simple flag transitivity]
\label{prop:simple}
For an abstract polyhedron, an automorphism fixing one flag is the
identity.  Hence a regular abstract polyhedron has a simply transitive
flag action.  If it is finite, then
\begin{equation}\label{eq:flags-group}
\abs{\Gamma(\Pcal)}=\abs{\Flag(\Pcal)}=4f_1,
\end{equation}
where \(f_1=\abs E\).
\end{proposition}

\begin{proof}
Rank preservation fixes every member of the chosen flag.  The diamond
condition then fixes each adjacent flag, and strong flag connectivity
propagates this to all flags and faces.  Thus a flag stabilizer is
trivial.  Each edge has two incident vertices and two incident faces,
so it belongs to four reduced flags.
\end{proof}

Fix a base flag.  For a regular \(\Pcal\), let \(\rho_i\) be the
unique automorphism sending it to its \(i\)-adjacent flag.  For type
\(\{p,q\}\),
\begin{equation}\label{eq:cgroup}
\rho_i^2=1,\qquad
(\rho_0\rho_2)^2=1,\qquad
\ord(\rho_0\rho_1)=p,\qquad
\ord(\rho_1\rho_2)=q,
\end{equation}
and the intersection property is
\begin{equation}\label{eq:intersection}
\langle\rho_i:i\in I\rangle\cap
\langle\rho_i:i\in J\rangle
=\langle\rho_i:i\in I\cap J\rangle .
\end{equation}
These data form a rank-\(3\) string C-group.  The orders \(p,q\) are
exact, and additional quotient relations may be present.  Therefore
\(\{p,q\}\) is local type data, not a complete global identifier.

\section{Abstract symmetry versus realization symmetry}

For a faithful nondegenerate realization \(\beta\), define
\[
G_\beta=
\{a\in\Isom(\R^3):a\text{ preserves the realized vertices, edges,
and declared faces}\}.
\]
Its action induces a subgroup
\(\iota_\beta(G_\beta)\leq\Aut(\Pcal)\).

\begin{definition}[Geometric regularity]
The realization \((\Pcal,\beta,\mathsf T)\) is geometrically regular
when \(G_\beta\) is transitive on its geometric flags.
\end{definition}

\begin{proposition}[One-way implication]
\label{prop:geometric-abstract}
A faithful geometrically regular realization has an abstractly
regular incidence polyhedron.  The converse fails.
\end{proposition}

\begin{proof}
The induced subgroup of \(\Aut(\Pcal)\) is already flag-transitive,
so the full automorphism group is flag-transitive.  Conversely, give
the abstract cube three pairwise distinct edge lengths in orthogonal
directions.  Its incidence automorphism group still has order \(48\)
and is flag-transitive, but the rectangular cuboid has only the eight
independent coordinate sign changes.  Edge directions cannot be
permuted, leaving six geometric flag orbits.
\end{proof}

Thus the legacy symbol \(G(\Pcal)\), described as ``geometric or
combinatorial depending on the setting'', is ill-typed.  The two
groups and the homomorphism between them must be named.

\section{Observation specifications and two preorders}

\begin{definition}[Typed observation specification]
An observation specification is
\[
\Ocal=(\Ccal_{\Ocal},Y_{\Ocal},Q_{\Ocal}),
\]
where \(\Ccal_{\Ocal}\) is a full subgroupoid of
\(\mathbf{RPol}_3\), stable under isomorphism, \(Y_{\Ocal}\) is a
declared data or label space, and
\[
Q_{\Ocal}:\Ob(\Ccal_{\Ocal})\longrightarrow Y_{\Ocal}
\]
is constant on the isomorphisms the protocol treats as nuisance
transformations.
\end{definition}

There are two independent comparisons.

\begin{align}
\Ocal_1\preceq_{\rm adm}\Ocal_2
&\quad\Longleftrightarrow\quad
\Ccal_{\Ocal_1}\text{ is a full subgroupoid of }\Ccal_{\Ocal_2},
\label{eq:admission}\\
\Ocal_{\rm fine}\succeq_{\rm info}\Ocal_{\rm coarse}
&\quad\Longleftrightarrow\quad
Q_{\rm coarse}=\kappa\circ Q_{\rm fine}
\text{ on a common domain}
\label{eq:information}
\end{align}
for some deterministic \(\kappa\).  Then
\(\Ker Q_{\rm fine}\subseteq\Ker Q_{\rm coarse}\): the coarse channel
identifies at least as many objects.

\begin{theorem}[Scoped filter monotonicity]
\label{thm:filter}
Let \(R\) be one fixed predicate, such as abstract regularity or
geometric regularity under one declared convention.  If
\(\Ocal_1\preceq_{\rm adm}\Ocal_2\), then
\[
\{X\in\Ob(\Ccal_{\Ocal_1}):R(X)\}
\subseteq
\{X\in\Ob(\Ccal_{\Ocal_2}):R(X)\}.
\]
No ordering of information kernels follows without
\cref{eq:information}, and no count comparison follows if different
quotient equivalences are used.
\end{theorem}

\begin{proof}
The set inclusion is immediate from inclusion of the object domains
while the predicate is held fixed.  The remaining claims concern
different data and therefore require a factorization map
\(\kappa\); domain inclusion alone supplies none.
\end{proof}

\begin{warning}[The historical family is not a chain]
Removing finiteness while retaining planarity and allowing
self-intersection while retaining finiteness relax different
predicates.  The resulting admissible classes may be incomparable.
The old ``filter ladder'' is therefore replaced by a filter poset.
\end{warning}

\begin{center}
\begin{tikzpicture}[
  node distance=7mm and 9mm,
  box/.style={draw=obsblue,rounded corners,fill=lightblue,
    align=center,inner sep=4pt,font=\scriptsize},
  arr/.style={-{Latex[length=2mm]},thick,draw=obsblue},
  dasharr/.style={-{Latex[length=2mm]},thick,dashed,draw=obsviolet},
  lab/.style={font=\tiny,fill=white,inner sep=1pt}
]
\node[box] (full) {full incidence\\and realization};
\node[box,right=of full] (skeleton) {vertex-edge\\skeleton};
\node[box,right=of skeleton] (image) {projected/noisy\\record};
\draw[arr] (full) -- node[above,lab] {forget faces} (skeleton);
\draw[arr] (skeleton) -- node[above,lab] {sensor channel} (image);
\node[box,below=of skeleton] (classes) {admissible object\\subgroupoid};
\draw[dasharr] (classes) -- node[right,lab] {domain choice} (skeleton);
\end{tikzpicture}
\end{center}

The horizontal arrows are lossy channels.  The vertical arrow is a
choice of domain.  Calling both operations a ``filter'' without types
caused the legacy order error.

\section{Classical spherical count and its hypotheses}

Let a finite equivelar polyhedral map of type \(\{p,q\}\) have
\(f\)-vector \((V,E,F)\).  Double counting incidences gives
\begin{equation}\label{eq:incidence}
pF=2E,\qquad qV=2E.
\end{equation}
If the underlying surface is the sphere,
\begin{equation}\label{eq:euler}
V-E+F=2,
\end{equation}
and therefore
\begin{equation}\label{eq:spherical}
\frac1p+\frac1q-\frac12=\frac1E>0.
\end{equation}
For integers \(p,q\ge3\), this leaves precisely
\[
(p,q)\in\{(3,3),(4,3),(3,4),(5,3),(3,5)\}.
\]
This proves the list of possible local types under the stated
spherical, finite, thin, equivelar hypotheses.  Existence and
uniqueness up to similarity of the corresponding convex regular
realizations are the separate classical realization theorem.
\Cref{eq:spherical} is not a proof classifying star polyhedra,
higher-genus maps, infinite tilings, or skeletal apeirohedra.

\begin{table}[ht]
\centering
\small
\caption{Exact finite controls for the five convex regular types.}
\begin{tabular}{lrrrrrr}
\toprule
Object & \(p\) & \(q\) & \(V\) & \(E\) & \(F\) & \(\abs{\Flag}=4E\)\\
\midrule
Tetrahedron & 3 & 3 & 4 & 6 & 4 & 24\\
Cube & 4 & 3 & 8 & 12 & 6 & 48\\
Octahedron & 3 & 4 & 6 & 12 & 8 & 48\\
Dodecahedron & 5 & 3 & 20 & 30 & 12 & 120\\
Icosahedron & 3 & 5 & 12 & 30 & 20 & 120\\
\bottomrule
\end{tabular}
\end{table}

Duality reverses the face order:
\[
\Pcal\longmapsto\Pcal^\ast,\qquad
(V,E,F)\longmapsto(F,E,V),\qquad
\{p,q\}\longmapsto\{q,p\},
\]
and preserves the automorphism-group order.  Duality is an exact
combinatorial operation, not an optical viewpoint change.

\section{Skeleton non-identifiability and the Petrial control}

For a regular map with distinguished generators
\((\rho_0,\rho_1,\rho_2)\), the Petrie operation replaces the
generating triple by
\begin{equation}\label{eq:petrie-generators}
(\rho_0\rho_2,\rho_1,\rho_2).
\end{equation}
Because \(\rho_0\) and \(\rho_2\) commute and are involutions,
\(\rho_0\rho_2\) is an involution.  On the map, the operation keeps
vertices and edges and selects Petrie polygons as faces.

\begin{proposition}[Explicit skeleton non-injectivity]
\label{prop:petrial}
Let \(Q_{\rm skel}(\Pcal)=(V,E)\).  The cube \(\Pcal_\square\) and its
Petrial \(\Pcal_\square^\pi\) satisfy
\[
Q_{\rm skel}(\Pcal_\square)
=Q_{\rm skel}(\Pcal_\square^\pi),
\qquad
\Pcal_\square\not\cong\Pcal_\square^\pi.
\]
The first has six square faces and Euler characteristic \(2\); the
second has four Petrie hexagons and Euler characteristic \(0\).
\end{proposition}

\begin{proof}
Both structures use the eight cube vertices and twelve cube edges.
The six square cycles give \(4\cdot6=24=2E\) edge-face incidences.
The four length-six Petrie cycles also give
\(6\cdot4=24=2E\).  Their \(f\)-vectors are respectively
\((8,12,6)\) and \((8,12,4)\), so they cannot be isomorphic as ranked
posets and have Euler characteristics \(2\) and \(0\).
\end{proof}

Thus ``same skeleton, different face observation'' is now an exact
counterexample, not a metaphor.  A two-dimensional rendering is even
coarser: projection can merge vertices, create apparent crossings,
or hide depth order.  Those effects belong to the sensor channel,
not to the incidence definition.

\section{Finite-resolution symmetry and identifiability}

Let observed vertices be
\[
y_i=x_i+\eta_i,\qquad \norm{\eta_i}\le\epsilon,
\]
after a declared registration, labeling, and incidence-matching
protocol.  For a candidate isometry \(g\) and allowed incidence
permutations \(\Pi\), define the RMS residual
\begin{equation}\label{eq:residual}
r_Y(g)=
\min_{\pi\in\Pi}
\left(\frac1n\sum_{i=1}^n
\norm{g y_i-y_{\pi(i)}}^2\right)^{1/2},
\qquad
\delta_Y(g)=\frac{r_Y(g)}{\ell_{\rm ref}},
\end{equation}
where \(\ell_{\rm ref}>0\) is declared.

\begin{proposition}[Exact symmetry under bounded noise]
\label{prop:noise}
If \(g x_i=x_{\pi(i)}\) for one allowed \(\pi\), then
\[
r_Y(g)\le2\epsilon.
\]
The converse is false: a small residual certifies only
threshold-dependent approximate symmetry.
\end{proposition}

\begin{proof}
Isometries preserve noise norms, so each residual vector is
\(g\eta_i-\eta_{\pi(i)}\) and has norm at most \(2\epsilon\).
The RMS bound follows.  Arbitrarily small asymmetric perturbations
give nonzero but arbitrarily small residuals.
\end{proof}

Let \(d_{\Qcal}\) be a declared metric on observations modulo nuisance
isometries and relabelings, and let
\(\{M_1,\ldots,M_K\}\) be a finite candidate library.  Define
\[
\Delta_i=\min_{j\ne i}d_{\Qcal}(M_i,M_j).
\]

\begin{theorem}[Half-separation certificate]
\label{thm:separation}
If \(\Delta_i>0\) and
\(d_{\Qcal}(Y,M_i)<\Delta_i/2\), then \(M_i\) is the unique nearest
candidate.
\end{theorem}

\begin{proof}
For \(j\ne i\), the triangle inequality gives
\[
d_{\Qcal}(Y,M_j)\ge
d_{\Qcal}(M_i,M_j)-d_{\Qcal}(Y,M_i)
>\Delta_i/2>d_{\Qcal}(Y,M_i).
\]
\end{proof}

This is conditional identifiability inside a declared finite library.
It does not show that the true object belongs to that library.

For a finite flag set and a declared group \(G\), let its flag orbits
have weights \(w_j\).  The diagnostic
\begin{equation}\label{eq:orbit-entropy}
H_{\rm orb}^{(2)}(G)
=-\sum_jw_j\log_2w_j
\end{equation}
vanishes exactly when there is one orbit.  It is a combinatorial
Shannon diagnostic, not thermodynamic entropy, and it is undefined
for an infinite flag set without a probability measure.

\section{Convention-scoped taxonomy and claim firewall}

\begin{table}[ht]
\centering
\small
\caption{Counts are meaningful only after the object convention is fixed.}
\begin{tabularx}{\linewidth}{L{30mm}L{43mm}Y}
\toprule
Count & Object convention & Permitted interpretation\\
\midrule
5 & finite convex regular polyhedra in \(\R^3\), up to similarity &
Classical Platonic classification.\\
4 & classical Kepler--Poinsot regular star polyhedra &
Star-face and self-intersection conventions are part of the type.\\
\(48=18+30\) & Schulte's discrete geometrically regular skeletal
polyhedra in \(\R^3\): 18 finite and 30 infinite &
Published classification in that specific skeletal convention.\\
Not 48 & all abstract regular rank-\(3\) polytopes or all
combinatorially regular realizations &
The abstract and realization universes are much larger.\\
\bottomrule
\end{tabularx}
\end{table}

\begin{warning}[Claim firewall]
This reference does not claim a new classification, a new regular
polyhedron, completeness beyond a cited convention, or recovery of
incidence from a picture.  Its strict internal results are the typed
filter preorders, the separation of abstract and geometric symmetry,
the explicit cube--Petrial non-injectivity control, and the
finite-resolution identification bounds.
\end{warning}

\section{Protocol minimum and benchmark audit}

A reproducible polyhedral observation claim must declare:
\begin{enumerate}[leftmargin=1.7em]
  \item the abstract object axioms and rank;
  \item the realization convention \(\mathsf T\), ambient space, and
  faithfulness;
  \item whether regularity uses \(\Aut(\Pcal)\) or \(G_\beta\);
  \item the admissible subgroupoid and the equivalence used for
  counting;
  \item the observation map, projection, labeling, and face-selection
  rule;
  \item spatial resolution, noise bound or law, registration group,
  estimator, threshold, and uncertainty;
  \item for infinite structures, discreteness, local finiteness,
  observation window, and boundary treatment.
\end{enumerate}

The executable companion checks all five Platonic incidence ledgers,
the spherical type inequality, dual pairs, all cube square and Petrie
cycles, the \(48\)-element cube action, the \(8\)-element generic
cuboid action, flag-orbit counts, bounded-noise residuals, and
half-separation certificates.  Passing these controls verifies the
claims encoded here; it is not an independent classification proof.

\section*{References}
\addcontentsline{toc}{section}{References}

{\raggedright\small
\begin{enumerate}[label={[\arabic*]},leftmargin=2em,itemsep=2pt,topsep=3pt]
  \item P. McMullen and E. Schulte,
  \emph{Abstract Regular Polytopes}, Encyclopedia of Mathematics and
  its Applications 92, Cambridge University Press, 2002.
  \href{https://doi.org/10.1017/CBO9780511546686}
  {doi:10.1017/CBO9780511546686}.

  \item P. McMullen and E. Schulte, ``Regular Polytopes in Ordinary
  Space'', \emph{Discrete \& Computational Geometry} 17 (1997),
  449--478.
  \href{https://doi.org/10.1007/PL00009304}
  {doi:10.1007/PL00009304}.

  \item E. Schulte, ``Regular Incidence Complexes, Polytopes, and
  C-Groups'', 2017.
  \href{https://arxiv.org/abs/1711.02297}{arXiv:1711.02297}.

  \item E. Schulte, ``Polyhedra, Complexes, Nets and Symmetry'',
  \emph{Acta Crystallographica A} 70 (2014), 203--216.
  \href{https://doi.org/10.1107/S2053273314000217}
  {doi:10.1107/S2053273314000217}.

  \item B. Grünbaum, ``Regular Polyhedra---Old and New'',
  \emph{Aequationes Mathematicae} 16 (1977), 1--20.
  \href{https://doi.org/10.1007/BF01836414}
  {doi:10.1007/BF01836414}.

  \item A. W. M. Dress, ``A Combinatorial Theory of Grünbaum's New
  Regular Polyhedra, Part I'', \emph{Aequationes Mathematicae} 23
  (1981), 252--265.
  \href{https://doi.org/10.1007/BF02188039}
  {doi:10.1007/BF02188039}.

  \item A. W. M. Dress, ``A Combinatorial Theory of Grünbaum's New
  Regular Polyhedra, Part II: Complete Enumeration'',
  \emph{Aequationes Mathematicae} 29 (1985), 222--243.
  \href{https://doi.org/10.1007/BF02189831}
  {doi:10.1007/BF02189831}.
\end{enumerate}
}

\end{document}
